\section{Discussion}
\label{sec:discussion}

HistoAtlas converts
routine H\&E pathology into 38 named, compartment-resolved measurements, attaches
a fully declared statistical layer to them, validates that layer in an
independent consortium, and publishes the result as queryable infrastructure with
provenance back to individual cells. The individual biological correspondences
reported here are, for the most part, expected.

What distinguishes this from another predictive computational-pathology system is
what the coordinates mean. Foundation-model embeddings are information-dense inputs for predictive
models~\cite{chen2024uni, vorontsov2024virchow}, but they are not interpretable:
their dimensions have no units,
cannot be compared between studies, and do not decompose into quantities a
pathologist can check, such as lymphocyte density inside tumor epithelium or
distance from a tumor cell to its nearest lymphocyte. Because HistoAtlas
measurements are named, they can be given confidence intervals that mean
something, corrected within families that can be audited, compared across
cohorts on their own scale, and traced back to the tissue that produced them. The cost is expressive
power: interpretable measurements will not match learned representations on pure
predictive benchmarks.

In the independent CPTAC cohort, the measurement distributions, the
correlations between measurements and the correlations with gene expression
were reproduced, two of six program distributions were unchanged, and
cluster-level pathway enrichment was not reproduced. At the measurement level, four of five matched cancer types passed the
prespecified concordance criterion and the inter-feature correlation geometry
was preserved, so the coordinates mean approximately the same thing on tissue
processed elsewhere. At the biological level, the strongest external evidence
came from molecular associations rather than from survival: all ten prespecified
histomic--transcriptomic associations replicated in direction and nine in
significance, effect sizes were concordant across 750 matched triples, and the
bridge extended to protein abundance with 98.8\% directional agreement among
doubly-significant pairs. Continuous programs transferred better than discrete
clusters, which is a useful negative result: it says that the reusable units of
this atlas are the measurements and their sparse combinations, not any
particular partition of feature space. Users should treat the ten-cluster
solution as a browsing convenience.

Several findings show what the layer is good for. The separation between
intratumoral and stromal lymphocyte density quantifies a compartment effect that
bulk tumor-infiltrating-lymphocyte scoring collapses~\cite{saltz2018til,
galon2006}; tumor--lymphocyte distance captures immune exclusion as a
geometric quantity rather than an impression; mitotic index tracks cell-cycle
transcripts and proteins; nuclear pleomorphism tracks luminal differentiation
markers, more closely at the protein level than at the transcript level, as
estrogen-receptor biology predicts. None of these is conceptually new. The
advance is that they are available systematically, pan-cancer, with units and
error bars, for anyone to query.

The proteomics analyses position the resource within the wider ecosystem. TCGA
anchored pan-cancer genomics and transcriptomics; CPTAC extended it to
proteomics~\cite{li2023cptac}; HistoAtlas adds morphology alongside them. The
finding that histomic measurements correlate coherently with protein abundance in
independent tissue argues that H\&E morphology is not merely a low-fidelity
shadow of RNA. But the stricter test constrains that claim sharply: after
adjustment for matched transcripts only 132 of \num{14687} feature--protein
triplets remained significant, concentrated in stromal and DNA-repair markers.
HistoAtlas is therefore not an alternative to transcriptomics or proteomics. It
is a widely available, spatially resolved, interpretable readout that tracks
many molecular states and adds a narrow increment beyond them. The survival
analyses point the same way: adding a measurement to a clinical-plus-RNA model
raised the apparent concordance index by a median of $+0.0007$ to $+0.0038$, with
gains in only about half of individual cancer types, and two of three prognostic
associations weakened toward the null once matched transcripts were added.

For researchers, the practical value is that a question requiring morphology can
now be asked without reprocessing images. A user can ask which measurements
associate with a given pathway in a given cancer type, compare a candidate
feature's behavior across 21 tumor types, check whether an association survives
subtype adjustment, or retrieve the minimum effect a given cohort could have
detected before concluding that a null result is informative. The atlas is also
available as a Model Context Protocol (MCP) server, which agentic systems can
query directly.

Several limitations bound what can be concluded. All associations are
observational and cross-sectional; none supports a causal claim. External
validation was substantially stronger for molecular associations than for
survival, because CPTAC follow-up was absent or minimal in the highest-priority
validation cohorts, so clinical-outcome claims rest primarily on TCGA. The
measurements inherit the errors of upstream segmentation: the cell model was
trained on eight cancer types, and 13 of the 21 included types were processed
out of distribution, with reduced detection sensitivity inflating kinetic indices
in some of them. Median tumor nuclear irregularity shifted by 1.8 to 2.4
standard deviations between cohorts and should not be compared across protocols.
Molecular subtype adjustment attenuated correlations by 28\% in pancreatic
adenocarcinoma against under 7\% in breast and endometrial carcinoma, so
subtype-aware reporting is needed in tumors with strong molecular substructure.
Cohort sizes are severely unequal, and 86\% of survival associations carry an
insufficient-evidence badge, reflecting power rather than absence of signal.
Molecular annotation is incomplete; analyzing one representative slide per
patient discards intra-patient heterogeneity, although the 312 TCGA patients
with several slides make it estimable in future releases; the nine-class taxonomy underlying the
measurements does not resolve macrophage or dendritic-cell populations that
matter for microenvironment biology, and the atlas has
not been benchmarked against foundation-model embeddings or against hybrid
representations that combine both.

Those limitations define the next steps: additional external cohorts to map
which measurements are protocol-robust, spatial transcriptomics as a stronger
ground truth for compartment-resolved immune and stromal quantities, prospective
and treatment-annotated cohorts to test clinical utility properly, and a direct
benchmark of interpretable measurements against learned embeddings to quantify
what interpretability costs. HistoAtlas as it stands is a reference layer: a
standardized, externally validated, openly queryable account of what routine
cancer pathology contains, built so that other people's questions can be asked
of it.
